\chapter{Simulation, Planning, and Model Error}
\label{ch:08}

The second family of instruments consists of models. A simulation is a computational representation of a system, constructed so that experiments which would be impossible, dangerous, or prohibitively slow in reality can be conducted on the representation, and planning is the use of such representations to choose among courses of action. The practice is ancient in spirit, since a map, a scale model, and a ship's log kept by dead reckoning are all instruments of this kind, and what is new is the scale, the speed, and the apparent authority of the computed result. The expectation that a system of high capability will transform planning is the reasoning behind many of the proposals in the later parts of this book, which are proposals for physical programs whose design depends on models of fluids, structures, ecosystems, orbits, and markets. The literature on decision making under deep uncertainty, which plans against many plausible models and not one~\cite{lempert2003}, supplies the discipline that the later chapters follow. It is therefore necessary to say clearly what a model can and cannot establish, and to state the discipline under which model-based plans should be accepted.

Every model has three sources of error that should be kept apart. The first is parametric error, the uncertainty in the numerical values of quantities that the model takes as input, such as the friction coefficient of a channel, the conductivity of a rock formation, or the strength of a material. Parametric uncertainty is the best understood, because it can be propagated through the model by sampling and its effect on the output can be reported as an interval. The second source is structural error, the error that arises because the equations of the model omit or misrepresent a process. A model of a river that neglects sediment transport may reproduce the discharge accurately and fail to predict that a channel will silt up, and no refinement of its parameters will repair the omission. Structural error is far more difficult to quantify, since by hypothesis the modeller has not thought of the omitted process, and the usual indicators of it are disagreement among structurally different models and failure of the model on cases that were not used to tune it. The third source is numerical error, introduced by the discretization and the finite precision of the computation, which can be assessed by refining the discretization and observing the convergence of the result.

For systems that exhibit sensitive dependence on initial conditions, a fourth consideration enters, which is the limit on predictability. If two states of a system that differ initially by a small amount $\delta_0$ separate at an exponential rate $\lambda$, then the error in a forecast grows until it reaches the scale $\Delta$ at which the forecast is no longer informative, and the time at which this occurs depends only logarithmically on the quality of the initial data. The consequence, familiar from the study of the atmosphere, is that improvements in the observation of the initial state buy additional forecasting time at a diminishing rate: a tenfold improvement in the initial data extends the horizon by the same fixed increment, whatever the starting accuracy. This is a property of the system and not of the model. It means that for some systems, no increase in computational capability can extend the reach of deterministic prediction beyond a horizon that can be calculated in advance, and that for these systems the proper objective of modelling is the prediction of statistics, such as the distribution of weather over a season or the frequency of floods in a century, and not of particular trajectories. Climate projection is a prediction of statistics in this sense, and its legitimacy does not depend on the predictability of weather, but the distinction has to be respected in the formulation of the questions that are put to models.

A model-based plan is accordingly a plan made in the presence of uncertainty about which model is correct, and the question of how to decide in these circumstances has a literature of its own. If the possible models can be assigned probabilities with which all parties agree, the plan that minimizes expected loss is the natural choice. But in the cases that matter most, the probabilities are exactly what is in dispute, and a plan that is optimal under one assignment can be disastrous under another. The decision criteria developed for such cases, grouped in the literature on decision-making under deep uncertainty, share the aim of selecting plans that perform acceptably across a wide range of plausible models and not plans that perform best under a single model. Among them is the criterion of minimax regret, which selects the plan that minimizes the largest shortfall, over models, between the plan's performance and the performance of the best plan for that model. It has the advantage of requiring no probabilities. It has the disadvantage of being sensitive to the set of models considered, since a model added to the set can change the result, and the set must therefore be constructed with care and with the participation of those who disagree about it. A further family of approaches seeks plans that are adaptive, meaning that they specify in advance the signals that will trigger revision of the plan and the actions that revision will take, so that the plan commits to a process of learning and not to a fixed trajectory.

Planning by a system of high capability does not remove these difficulties and in one respect aggravates them. A system that can explore a very large space of plans will find those that perform best under the model it has been given, and plans that perform best under a model tend to be those that exploit its errors. The phenomenon, known in several fields as overfitting to the simulator or as the exploitation of model artifacts, is the consequence of optimization pressure applied to an imperfect representation. A control strategy found by optimization against a simulation of a bridge may be brittle in the real structure for reasons that no one anticipated, and a design found by optimization against a hydrological model may rely on a numerical feature of the model that has no counterpart in nature. The protection is not to avoid optimization but to subject the plans that optimization finds to tests the optimizer did not see: models of different structure, physical experiments at reduced scale, and staged deployment in which the plan is first applied to a small part of the system and its performance is compared with the prediction. The discipline is stated in Chapter 10 as a bound on the fidelity gap, and it is reflected in the readiness classification of Chapter 16, in which projects are not permitted to advance until their central models have been tested against measurements.

A last observation concerns the social role of models. A model embodies assumptions, and when it is used to justify a decision the assumptions become a hidden part of the decision. The citizen who is told that a plan was chosen because a computation showed it to be optimal is in no position to contest the plan unless the assumptions can be stated, the data exposed, and the sensitivity of the result to the assumptions shown. The requirement that a model-based recommendation be accompanied by this information is a requirement of accountability in the sense of Chapter 2: the party that adopts the recommendation remains answerable for the decision, and answerability requires access to the grounds. The advance of capability in modelling should therefore be accompanied by an advance in the practice of documentation, in which models intended for public decisions are versioned, archived, and open to independent execution. A model that cannot be re-run by a skeptical party is an oracle, and the program of this book assigns no authority to oracles.

\section*{Formalization}

Two results are stated, the first on the limit of predictability and the second on decision under model uncertainty.

\begin{definition}[Predictability horizon]
Let the separation of two trajectories obey $\delta(t)=\delta_0e^{\lambda t}$ with $\lambda>0$ until it reaches the saturation scale $\Delta>\delta_0$. The \emph{predictability horizon} is the time $T(\delta_0)$ at which $\delta(T)=\Delta$.
\end{definition}

\begin{proposition}
$T(\delta_0)=\lambda^{-1}\ln(\Delta/\delta_0)$. For any $r>1$, reducing the initial error from $\delta_0$ to $\delta_0/r$ increases the horizon by exactly $\lambda^{-1}\ln r$, independent of $\delta_0$. To extend the horizon by a prescribed amount $\tau$ the initial error must be reduced by the factor $e^{\lambda\tau}$.
\end{proposition}

\begin{proof}
Setting $\delta_0e^{\lambda T}=\Delta$ gives the formula, and $T(\delta_0/r)-T(\delta_0)=\lambda^{-1}\ln r$ follows by subtraction. The last statement is the case $\ln r=\lambda\tau$.
\end{proof}

With $\lambda=0.5\ \mathrm{day^{-1}}$, a figure of the order of magnitude associated with the error doubling of mid-latitude weather over a few days, a tenfold reduction in initial error extends the horizon by $\ln10/0.5\approx4.6$ days, and a hundredfold reduction by $9.2$ days. The cost of observation, which increases at least in proportion to the reduction achieved, therefore buys additional forecast time at a rate that falls exponentially, and beyond a horizon set by the largest feasible reduction the forecast of individual trajectories is unavailable and only statistics remain.

For decision under uncertainty, let $\mathcal A$ be a finite set of actions and $\mathcal M$ a finite set of models, with loss $L(a,m)\ge0$. Define the \emph{regret} $r(a,m)=L(a,m)-\min_{a'}L(a',m)$ and the \emph{minimax-regret action} $a^\ast\in\arg\min_a\max_mr(a,m)$ with value $r^\ast=\min_a\max_mr(a,m)$.

\begin{proposition}
For every probability distribution $\pi$ on $\mathcal M$, the Bayes action $a_\pi\in\arg\min_a\E_\pi L(a,m)$ satisfies $\E_\pi r(a_\pi,m)\le r^\ast$. Hence the expected regret of a decision maker who chooses by any prior is bounded by the minimax regret, which can be reported as the price of declining to commit to a prior.
\end{proposition}

\begin{proof}
Since $\E_\pi r(a,m)=\E_\pi L(a,m)-\E_\pi\min_{a'}L(a',m)$, the Bayes action for loss also minimizes expected regret. Therefore $\E_\pi r(a_\pi,m)\le\E_\pi r(a^\ast,m)\le\max_mr(a^\ast,m)=r^\ast$.
\end{proof}

A numerical illustration shows that the criteria can disagree. Let there be two models and three actions with losses $L(a_1,\cdot)=(0,6)$, $L(a_2,\cdot)=(4,0)$, $L(a_3,\cdot)=(2,3)$. The regrets are $(0,6)$, $(4,0)$, $(2,3)$, so that the maximum regrets are $6,4,3$ and the minimax-regret action is $a_3$ with $r^\ast=3$. Under equal weights the expected losses are $3$, $2$, and $2.5$, so that the Bayes action is $a_2$, whose maximum regret is $4$ and whose expected regret is $2$, below $r^\ast$ as the proposition requires. A decision maker unwilling to commit to equal weights, and who would regret a loss of $4$ under the first model more than the cost of insurance, would choose $a_3$. The example illustrates that the selection of a criterion is a normative decision which belongs to those who bear the consequences and not to the modeller, a conclusion that supports the allocation of authority described in Chapter 13.
